\chapter{การตรวจวัดและการวิเคราะห์สเปกตรัมภาพถ่ายเกษตรแม่นยำ}
\label{chap:ch03}

% --- หน้าที่ 1 ของบท: แผนบริหารการสอน (กล่อง 1 และ 2) ---
\section*{แผนบริหารการสอนประจำบทที่ 3}
\addcontentsline{toc}{section}{แผนบริหารการสอนประจำบทที่ 3}

\begin{planbox}{หัวข้อเนื้อหาประจำบท}
\begin{enumerate}
    \item ทฤษฎีการสะท้อนแสงของพืชพรรณและแถบสเปกตรัมที่ตามองเห็น
    \item ดัชนีความต้านทานบรรยากาศ VARI สำหรับภาพถ่ายการเกษตรมุมสูง
    \item ดัชนีใบเขียว GLI และดัชนีความเขียวส่วนเกิน ExG
    \item ดัชนีความสมบูรณ์พืช Visible NDVI และการตัดแยกทรงพุ่ม Canopy Mask
    \item สถาปัตยกรรมประมวลผลภาพคู่ขนาน Pure Dart Isolate Computing
\end{enumerate}
\end{planbox}

\begin{planbox}{วัตถุประสงค์เชิงพฤติกรรม}
\begin{enumerate}
    \item อธิบายพฤติกรรมการดูดกลืนและการสะท้อนแสงของสารคลอโรฟิลล์ในใบพืชได้อย่างถูกต้อง
    \item คำนวณค่าดัชนีสเปกตรัม VARI, GLI, ExG และ Visible NDVI จากค่าพิกเซลได้อย่างแม่นยำ
    \item สร้างภาพแผนที่ความร้อน False-Color Heatmap เพื่อจำแนกสุขภาพพืชได้อย่างสมบูรณ์
    \item พัฒนาระบบประมวลผลภาพแบบ Isolate ในภาษา Dart เพื่อป้องกันการกระตุกของหน้าจอได้อย่างมีประสิทธิภาพ
\end{enumerate}
\end{planbox}

\newpage

% --- หน้าที่ 2 ของบท: แผนบริหารการสอน (กล่อง 3, 4, 5) ---
\begin{planbox}{วิธีสอนและกิจกรรมการเรียนการสอน}
\begin{enumerate}
    \item การบรรยายหลักการทางทัศนศาสตร์ของการสำรวจระยะไกลเชิงเกษตรกรรม
    \item การทดลองวิเคราะห์ภาพถ่ายแปลงทุเรียนตัวอย่างด้วยอัลกอริทึมสเปกตรัมชนิดต่าง ๆ
    \item การฝึกปฏิบัติการปรับแต่งค่าความไว Threshold เพื่อคัดกรองผิวดินออกจากทรงพุ่มพืช
    \item การอภิปรายเปรียบเทียบข้อดีและข้อจำกัดของดัชนีพืชพรรณแต่ละชนิดในแปลงจริง
\end{enumerate}
\end{planbox}

\begin{planbox}{สื่อการเรียนการสอน}
\begin{enumerate}
    \item ชุดภาพถ่ายทางอากาศแปลงเกษตรกรรมความละเอียดสูงในย่านแสงสีแดง เขียว และน้ำเงิน
    \item โมดูลซอฟต์แวร์ AgriVisionService สำหรับการประมวลผลภาพบน Flutter
    \item สื่อภาพจำลองการดูดกลืนแสงของสารสีในพืชและกราฟคุณลักษณะสเปกตรัม
    \item เอกสารประกอบการสอนเรื่องดัชนีพืชพรรณในย่านแสงที่ตามองเห็น
\end{enumerate}
\end{planbox}

\begin{planbox}{การวัดและประเมินผล}
\begin{enumerate}
    \item การประเมินผลการคำนวณดัชนีสเปกตรัมจากชุดข้อมูลตัวเลขในแบบฝึกหัด
    \item การประเมินความถูกต้องในการแปลผลระดับความสมบูรณ์ของพืชจากแผนที่ความร้อน
    \item การประเมินความกระตือรือร้นในการมีส่วนร่วมและการตอบคำถามท้ายบทเรียน
\end{enumerate}
\end{planbox}

\clearpage

% --- หน้าที่ 3 เป็นต้นไป: เนื้อหาวิชาการ ---
\section{ทฤษฎีการสะท้อนแสงของพืชพรรณและแถบสเปกตรัม}

หลักการทางฟิสิกส์ของการสำรวจระยะไกลด้วยภาพถ่ายทางอากาศ อาศัยความสัมพันธ์ระหว่างการดูดกลืนและการสะท้อนของคลื่นแม่เหล็กไฟฟ้าต่อโครงสร้างทางชีววิทยาของใบพืช สารคลอโรฟิลล์ (Chlorophyll a และ Chlorophyll b) ภายในเซลล์พืชจะดูดกลืนพลังงานแสงในย่านสีน้ำเงิน (ความยาวคลื่นประมาณ 450 นาโนเมตร) และย่านสีแดง (ความยาวคลื่นประมาณ 660 นาโนเมตร) อย่างรุนแรงเพื่อใช้ในกระบวนการสังเคราะห์ด้วยแสง ในขณะที่จะสะท้อนพลังงานแสงในย่านสีเขียว (ความยาวคลื่นประมาณ 550 นาโนเมตร) ออกมามากกว่า จึงทำให้สายตามนุษย์มองเห็นใบพืชเป็นสีเขียว

เมื่อพืชเกิดสภาวะเครียด (Plant Stress) อันเนื่องมาจากการขาดน้ำ ขาดธาตุอาหารไนโตรเจน หรือการถูกทำลายโดยโรคและแมลง โครงสร้างเซลล์และปริมาณคลอโรฟิลล์จะลดลง ส่งผลให้อัตราการดูดกลืนแสงสีแดงลดลงและสะท้อนแสงสีแดงออกมามากขึ้น ทำให้ใบเปลี่ยนเป็นสีเหลืองหรือสีน้ำตาล

\begin{figure}[H]
\centering
\resizebox{0.86\textwidth}{!}{
\begin{tikzpicture}[>=Stealth]
    % Axes
    \draw[->, line width=1.2pt] (0,0) -- (10,0) node[right] {\footnotesize ความยาวคลื่น (nm)};
    \draw[->, line width=1.2pt] (0,0) -- (0,5.5) node[above] {\footnotesize การสะท้อนแสง (\%)};
    
    % Bands Regions
    \fill[blue!15] (1,0) rectangle (3,5);
    \node at (2,4.7) {\footnotesize น้ำเงิน (Blue)};
    \fill[green!20] (3,0) rectangle (6,5);
    \node at (4.5,4.7) {\footnotesize เขียว (Green)};
    \fill[red!15] (6,0) rectangle (9,5);
    \node at (7.5,4.7) {\footnotesize แดง (Red)};

    % Healthy Plant Curve
    \draw[line width=2pt, color=agriGreen] plot [smooth] coordinates {
        (0.8,0.8) (2,0.6) (4.5,3.2) (7,0.8) (9,2.8) (9.8,5.0)
    };
    \node[color=agriGreen, font=\footnotesize\bfseries] at (5.2,3.6) {พืชสมบูรณ์ (Healthy)};

    % Stressed Plant Curve
    \draw[line width=2pt, dashed, color=agriYellow!80!black] plot [smooth] coordinates {
        (0.8,1.2) (2,1.2) (4.5,2.4) (7,2.2) (9,2.6) (9.8,3.5)
    };
    \node[color=agriYellow!80!black, font=\footnotesize\bfseries] at (7.5,2.6) {พืชเครียดขาดน้ำ (Stressed)};
\end{tikzpicture}
}
\caption{กราฟลักษณะการสะท้อนแสงของพืชสมบูรณ์เปรียบเทียบกับพืชในสภาวะเครียด}
\label{fig:spectral_reflectance}
\end{figure}

\section{ดัชนีความต้านทานบรรยากาศสำหรับภาพถ่ายมุมสูง}

เนื่องจากกล้องที่ติดตั้งบนโดรนขนาดเล็กส่วนใหญ่เป็นเซนเซอร์แถบแสงที่ตามองเห็น (RGB Sensor) ระบบจึงใช้ดัชนีพืชพรรณในย่านแสงที่ตามองเห็น (Visible Vegetation Indices) โดยดัชนีที่ได้รับการยอมรับและมีประสิทธิภาพสูงสุดในการบินระดับความสูงคือ \textbf{Visual Atmospheric Resistance Index (หรือ VARI)}

\begin{formulabox}[สมการดัชนีความต้านทานบรรยากาศ VARI]
\begin{equation}
\text{VARI} = \frac{\text{Green} - \text{Red}}{\text{Green} + \text{Red} - \text{Blue}}
\label{eq:vari}
\end{equation}
โดยที่ $\text{Green}, \text{Red}, \text{Blue}$ คือค่าความสว่างของจุดภาพในแถบแสงสีเขียว สีแดง และสีน้ำเงิน
\end{formulabox}

จุดเด่นของดัชนี VARI คือการนำค่าในแถบแสงสีน้ำเงินมาเป็นพจน์ปรับเทียบในตัวหาร ซึ่งช่วยหักล้างการกระเจิงของแสงจากฝุ่นละออง ไอน้ำ และหมอกควันในชั้นบรรยากาศ (Atmospheric Scattering) ได้อย่างดีเยี่ยม ทำให้ผลการวิเคราะห์มีความเสถียรแม้ในสภาพอากาศที่มีแดดจัดหรือมีหมอกควันปกคลุม

\section{ดัชนีใบเขียวและดัชนีความเขียวส่วนเกิน}

นอกจากดัชนี VARI แล้ว ระบบยังรองรับดัชนีเสริมอื่นที่มีความสำคัญในการจำแนกพืชพรรณ ได้แก่

\begin{enumerate}
    \item \textbf{Green Leaf Index (หรือ GLI)} ออกแบบมาเพื่อเน้นการคัดแยกใบไม้สีเขียวออกจากเศษซากพืชแห้งและผิวดิน
    \begin{formulabox}[สมการดัชนีใบเขียว GLI]
    \begin{equation}
    \text{GLI} = \frac{2\cdot\text{Green} - \text{Red} - \text{Blue}}{2\cdot\text{Green} + \text{Red} + \text{Blue}}
    \label{eq:gli}
    \end{equation}
    \end{formulabox}

    \item \textbf{Excess Green Index (หรือ ExG)} เป็นดัชนีที่เหมาะสำหรับการทำภาพแยกส่วน (Binarization Segmentation) เพื่อแยกแยะทรงพุ่มออกจากพื้นดิน
    \begin{formulabox}[สมการดัชนีความเขียวส่วนเกิน ExG]
    \begin{equation}
    \text{ExG} = 2\cdot\text{Green} - \text{Red} - \text{Blue}
    \label{eq:exg}
    \end{equation}
    \end{formulabox}
\end{enumerate}

\begin{table}[H]
\caption{การเปรียบเทียบคุณสมบัติและการใช้งานของดัชนีพืชพรรณในย่านแสงที่ตามองเห็น}
\label{tab:indices_comparison}
\begin{tabularx}{\textwidth}{p{2.8cm}p{3.8cm}p{3.5cm}Y}
\toprule
\textbf{ชื่อดัชนี} & \textbf{ช่วงค่าตัวเลข} & \textbf{จุดเด่นสำคัญ} & \textbf{การนำไปประยุกต์ใช้งาน} \\
\midrule
VARI & $-1.0$ ถึง $+1.0$ & ทนต่อการรบกวนของชั้นบรรยากาศ & วิเคราะห์ความสมบูรณ์และคลอโรฟิลล์ \\
GLI & $-1.0$ ถึง $+1.0$ & แยกใบสดออกจากเศษฟางและวัชพืชแห้ง & ตรวจสอบการเติบโตของพืชแปลงปลูก \\
ExG & $-255$ ถึง $+510$ & ประมวลผลรวดเร็ว แยกทรงพุ่มชัดเจน & สร้าง Canopy Mask วัดพื้นที่ปกคลุม \\
Visible NDVI & $-1.0$ ถึง $+1.0$ & มีโครงสร้างใกล้เคียงกับดัชนี NDVI & ประเมินสัดส่วนความเขียวขจีรวม \\
\bottomrule
\end{tabularx}
\end{table}

\section{การตัดแยกทรงพุ่มและการแปลงเป็นแผนที่ความร้อน}

การประเมินสัดส่วนการปกคลุมดินของทรงพุ่ม (Canopy Coverage Percentage) อาศัยการสร้างหน้ากาก \textbf{Canopy Mask} โดยหากจุดภาพใดมีค่าดัชนีสเปกตรัมสูงกว่าค่าความไวขีดเริ่มเปลี่ยน (Sensitivity Threshold) จะถูกจัดเป็นจุดภาพของทรงพุ่มพืช ส่วนจุดภาพที่มีค่าต่ำกว่าจะถูกจัดเป็นผิวดินหรือเศษหิน

จากนั้นระบบจะแปลงค่าดัชนีของแต่ละจุดภาพให้เป็นสีแบบจำลอง \textbf{False-Color Palette} เพื่อให้นักบินมองเห็นความแปรปรวนในแปลงได้อย่างชัดเจน โดยไล่ระดับสีจากสีเขียวเข้ม (สมบูรณ์สูง คลอโรฟิลล์หนาแน่น) ผ่านสีเขียวอ่อน สีเหลือง และสีแดง (พืชเกิดสภาวะเครียด ขาดน้ำ หรือใบเหลือง)

\section{สถาปัตยกรรมประมวลผลภาพคู่ขนานแบบไอโซเลต}

เนื่องจากการคำนวณดัชนีสเปกตรัมสำหรับภาพถ่ายขนาด 1 ล้านพิกเซลขึ้นไป ต้องใช้การดำเนินการทางคณิตศาสตร์นับล้านครั้ง หากรันบนเธรดหลัก (UI Thread) จะส่งผลให้หน้าจอควบคุมการบินเกิดอาการค้างหรือกระตุก ซึ่งเป็นอันตรายอย่างยิ่งต่อการควบคุมโดรน

แอปพลิเคชันจึงแก้ปัญหานี้ด้วยการใช้กลไก \textbf{Flutter Compute Isolate} ซึ่งแยกงานประมวลผลอาร์เรย์พิกเซลไบต์ \texttt{Uint8List} ไปรันบนหน่วยประมวลผลย่อยในเบื้องหลังแบบคู่ขนาน ทำให้หน้าจอ HUD ยังคงสามารถอัปเดตการแสดงผลและส่งแพ็กเก็ตควบคุมการบินได้ต่อเนื่องอย่างราบรื่นที่ 60 เฟรมต่อวินาที

\clearpage

\section*{คำถามทบทวนท้ายบทที่ 3}
\addcontentsline{toc}{section}{คำถามทบทวนท้ายบทที่ 3}

\begin{enumerate}
    \item จงอธิบายบทบาทของคลอโรฟิลล์ในการดูดกลืนแสงสีแดงและสะท้อนแสงสีเขียว พร้อมวาดภาพประกอบโดยสังเขป
    \item จุดภาพพิกเซลหนึ่งมีค่า RGB เท่ากับ (80, 160, 60) จงคำนวณค่า VARI, GLI และ ExG ของจุดภาพนี้ พร้อมระบุว่าพืชอยู่ในเกณฑ์ใด
    \item เหตุใดดัชนี VARI จึงมีความทนทานต่อการรบกวนของฝุ่นละอองในชั้นบรรยากาศมากกว่าดัชนี ExG
    \item จงอธิบายความสำคัญของการใช้ Isolate ในการประมวลผลภาพถ่ายการเกษตรในระบบ Flutter
\end{enumerate}

\clearpage
